\documentclass[11pt]{article}
\usepackage[margin=1in]{geometry}
\usepackage{amsmath,amssymb}
\usepackage{booktabs}
\usepackage{xcolor}
\usepackage[colorlinks=true,linkcolor=black,urlcolor=blue]{hyperref}
\usepackage{parskip}
\newcommand{\xor}{\oplus}
\title{\textbf{Nonlinear gadgetization}\\[2pt]
\large The product-share encoding, described from scratch}
\author{}
\date{2026-07-24}
\begin{document}
\maketitle

\noindent\emph{Written for a reader who knows the linear (XOR-pair) gadgetizer
and nothing else about this line of work. Covers the design of 2026-07-23 and
the revisions of 2026-07-24 as one piece. Option reference in \S9.}

\section{Where linear gadgetization leaks}

The gadgetizer takes a small circuit $C$ and produces a much larger equivalent
$G$ that computes the same function on its low wires. Every logical value of $C$
is \textbf{secret-shared} across the wires of $G$, and $G$'s gates manipulate the
shares. The security goal is that $G$ reveal nothing about $C$ beyond its
input/output behaviour---in particular that an adversary looking at a point of
$G$ cannot tell \textbf{where in $C$'s computation} that point sits.

In the linear gadgetizer a logical value $v$ lives on two carrier wires as
$v=c_0\xor c_1$. This is the whole problem. The decode is linear, so
\textbf{at every instant, every logical value of $C$ is a degree-1 (affine)
function of $G$'s wires.} An adversary who fits affine functions can therefore
read $C$'s intermediate state out of $G$ at any point---no matter how
algebraically complex that state is as a function of the input.

\texttt{hmap\_affine} measures this directly. For each prefix $i$ of $C$ and
prefix $j$ of $G$, it fits the best $\mathrm{GF}(2)$ affine reconstruction of
$C$'s state after $i$ gates from $G$'s wire values after $j$ gates, and reports
held-out error $H$ ($0$ = perfectly recoverable, $0.5$ = hidden). Under linear
sharing the aligned cells light up and a low-$H$ \textbf{diagonal} runs across
the map. That diagonal is a picture of $C$'s computational progress inside $G$.

Read these maps by the \textbf{ridge}, never by the mean (which saturates near
$0.5$): \emph{depth} is the per-row prominence of the low-$H$ valley above its
row background, $\rho$ is the Spearman correlation of row index against ridge
location (``is it a monotone diagonal''), and \emph{perm-z} scores that against
a shuffled null.

\textbf{The diagonal survived everything.} Local mixing (ssg, fmix), generation
targeting to generation 100+, four million mixing moves, growth phases: all of
them shallow the valley somewhat but leave $\rho=1$. The reason is
structural---mixing re-randomizes \emph{which} wires carry the sharing, but the
value stays affine in the new wires, so an affine adversary just re-fits.
Masking values \emph{between} their uses does not help either: the ridge lives
\textbf{at} the use points, and any scheme that un-masks an operand so an
ordinary gate can read it re-exposes exactly there. (An earlier ``deferred
mask'' design failed on precisely this.)

The leak is not a mixing deficiency. It is the linear decode.

\section{The fix: a balanced nonlinear decode}

Give each logical value, on top of its XOR pair, a set of permanent
\textbf{multiplicative mask terms}:
\[
v_i \;=\; w_u \xor w_{u'} \;\xor\; \bigoplus_j \prod_l \bigl(w_{p_{jl}} \xor a_{jl}\bigr) \;\xor\; c_i
\]
Each $\prod_l (w\xor a)$ is a conjunction of literals over a dedicated band of
wires; $c_i$ is a bookkeeping constant held in a ledger, never on a wire. Now
the decode is nonlinear, and an affine adversary cannot express it.

\textbf{Why the XOR pair is still there.} A pure product decode is un-gateable,
and the reason is a counting obstruction worth stating: $(w_p\xor a)(w_q\xor b)$
partitions its representations into classes of size 3 and 1, and no reversible
gadget can conditionally flip between classes of \emph{unequal} size---a
bijection cannot map a 3-element class onto a 1-element class. This holds for
every fixed $(a,b)$, with any ancillae, garbage or re-encoding. The linear part
re-symmetrizes the classes; it is structural, not a convenience.

\section{The frozen source band}

The mask sources live on a band of extra wires appended after the $2n$ carriers.
Each band wire is filled at the input port and then \textbf{never written
again}. Because the sources are frozen, every registered mask term is
\textbf{time-invariant}---unaffected by all re-randomization and all gate
traffic---so a mask is one gate, not a maintained data structure, and nothing is
ever flushed except the final unshare.

\textbf{The fill must not be linear.} The obvious fill, a random XOR of data
wires, makes every band wire an \emph{affine invariant of the entire circuit},
held from the input port to the output port. That is exactly the class of
relation which \textbf{collapses the preimage SAT search}: an attacker who
learns such invariants from forward runs and injects them turned $n{=}32$
instances from ``$>$2 hours, timeout'' into ``$\sim$10 minutes''. A linear band
fill manufactures them for free.

The fill is therefore
\[
\mathit{band}_j \;=\; \mathit{junk}\;\xor\;x_{\mathrm{pivot}}\;\xor\;(\text{small linear part})\;\xor\;\bigoplus\nolimits^{M}(\text{two-source products})
\]
with each product's two sources drawn from the data wires \textbf{and from
already-filled band wires}. The cascade is the point: multiplying two
already-balanced band bits keeps the firing rate near $1/4$ while the input
degree multiplies up the band, so high degree in $x$ is reached with nothing
wider than a two-control gate. (A flat high-degree monomial would fire on only a
$2^{-d}$ fraction of inputs and behave statistically like its linear part.)

Balance is exact, which the mask argument needs: the fresh pivot is excluded
from the rest of that wire's transitive support, so the fill is balanced over
uniform $x$ for any nonlinear part.

The same fill is emitted again after the output bookend (\textbf{mirror $F'$}),
so the band is not anchored only at one port.

\section{The share-native fold---the part that actually matters}

An encoding is only as good as the gate that operates on it. If a gate
reconstructs its operands to compute with them, the ridge comes straight back at
the use point.

So the gadget never reconstructs. A source gate $v_t \mathrel{\xor}= f(\text{controls})$
is applied by \textbf{expanding $f$'s ANF over the operands' full decodes}. Each
control value contributes a set of atoms---its two carrier literals, its mask
terms, and a constant when the ledger constant or a negative control polarity
applies---and every cross-product of one atom per control becomes \textbf{one
conjunction fragment} written into a randomly chosen carrier of the target.

The property that matters: a value's own mask term appears only \emph{inside} a
fragment's firing condition, as control literals. \textbf{It is never
materialized on a wire.} No fragment computes a value's two-carrier XOR, and no
wire is written except the target's carriers. Operands stay masked straight
through the gate. Constants and the gate's complement go to the ledger constant,
so the encoding emits no bare $X$ anywhere.

Fragments are emitted in randomized order. Without this, each source gate
produced its $(k+2)^2$ fragments consecutively, in deterministic order, all
aimed at one value's two carriers---a \textbf{progress clock readable
statically}, with no execution required.

\section{Degree and count do different jobs}

Two independent levers, and conflating them wastes gates.

\textbf{Degree hides algebraically.} A degree-$d$ mask term is not in the
degree-$<d$ $\mathrm{GF}(2)$ span, so a \emph{single} one forces $H=0.5$ against
\textbf{every} reconstruction adversary of degree $<d$, regardless of how many
terms there are.

\textbf{Count hides statistically.} Against a best-approximation adversary, $k$
stacked terms pile up: the best affine readout error tends to $1/2$ as
$1/2-2^{-(k+1)}$.

This yields \textbf{towers}: a degree-$d$ mask kills every adversary below
degree $d$, and the diagonal returns at degree exactly $d$. The wall moves up
one level per $+1$ of mask degree; it does not vanish. Verified with real
degree-2 and degree-3 adversaries, each capability-controlled.

A degree-$d$ term is also a \textbf{sparser} perturbation (it flips the value
only $2^{-d}$ of the time), so wider masks are individually \emph{weaker}
statistically even as they are stronger algebraically. Hence a \textbf{mixed}
design: a few low-degree terms for statistical margin plus tower terms for
algebraic hiding. The production plan is $[2,3,3]$---one degree-2 base term and
two degree-3 tower terms. Two degree-3 terms are what kill the degree-2
diagonal; one leaves a boundary residual.

Raising the tower is nearly free in gate count (fragment \emph{count} is set by
the term count and the source gate's arity, not by mask degree---widening a mask
only adds literals to existing fragments). What it costs is fragment
\textbf{width}.

\section{Gate width, and a corrected assumption}

Folding a degree-$d$ mask through a $w$-operand gate yields fragments of width up
to $d\cdot w$, so the production gadget contains gates of width up to 6. Two
consequences were expected to be bad, and exactly one of them is.

\textbf{Real:} with \texttt{--db-ctrl-cap 2}, any gate with $\ge3$ controls is
ineligible for frozen-DB re-encoding. About 62\% of the gadget is in that class.
This is why generation targeting cannot report progress here (\S8)---and, in the
other direction, why the encoding \textbf{resists DB-based simplification},
which is a property one wants of an obfuscation.

\textbf{Not real:} the assumption that wide gates are immobile fossils. A
per-width commutation-leeway census of the real $n{=}128$ gadget says the
opposite---the width-5 and width-6 gates are the \textbf{most mobile gates in
the circuit} (mean float-box 2779 and 2778, versus 694 for the g57s). A gate
with many literals has more chances to hold an \emph{opposite} literal against a
neighbour, which exempts the pair from colliding, so it crosses more freely.

A mode exists that removes wide gates entirely (\texttt{--prod-max-width}, \S9):
it ladders every emission down to two controls using a Barenco double sweep over
\textbf{dirty borrowed carriers}---arbitrary-valued wires, each visited an even
number of times so its dirty content cancels between readings. It costs no extra
wires and preserves correctness under arbitrary junk. But it costs
$\mathbf{6.4\times}$ the gates and \emph{reduces} mobility to mean 259, so it is
off by default. Its one live justification is DB re-encodability, and whether
that is worth $6.4\times$ is not yet measured.

\section{What it achieves}

All heatmaps read by the ridge measure with verified references, on a real
$n{=}128$ sliced sandwich (source $C$ = 3000 g57 gates).

\begin{center}
\begin{tabular}{llrrrr}
\toprule
gadget & adversary & depth & depthMed & $\rho$ & meanH \\
\midrule
linear (plain), 85{,}169 g / 512 w & degree 1 & 0.346 & 0.356 & \textbf{1.000} & 0.475 \\
$[2,3,3]$, 281{,}068 g / 568 w & degree 1 & 0.044 & \textbf{0.000} & 0.015 & 0.499 \\
$[2,3,3]$, 281{,}068 g / 568 w & degree 2 & 0.043 & \textbf{0.000} & $-0.007$ & 0.498 \\
\bottomrule
\end{tabular}
\end{center}

The median original-circuit prefix has \textbf{zero} ridge prominence and the
reconstruction is uncorrelated with progress. The nulls are
capability-controlled: a design whose masks are only degree 2 \emph{is} recovered
by the same degree-2 adversary at $\rho=1.00$, so these are genuine hiding, not
a weak measurement.

Earlier work carried the whole pipeline through---gadget $\to$ a DB-mixing phase
$\to$ threefold growth---and the diagonal stayed dead \textbf{identically} at
every stage. Mixing neither helps nor hurts, because the leak is gone before
mixing begins.

\textbf{Cost.} The encoding multiplies the plain gadget by ${\sim}3.3\times$ in
gates and adds a band of extra wires (56 at $n{=}128$, versus 512 carriers). Per
source gate the fold emits about $(k_{\mathrm{total}}+2)^2$ fragments, so the
term count is a quadratic lever and everything else is second order.

\section{Reading a mixing run on this material}

\texttt{G=} in the fmix report is the 5th-percentile gate generation over
\textbf{all} gates. Since ${\sim}62\%$ of a product-share gadget is
cap-ineligible and never re-encoded, \texttt{G=} is \textbf{pinned at 0 by
construction} and no amount of mixing moves it. ``Run to generation $N$'' cannot
mean the reported \texttt{G=} here.

The meaningful target is generation $N$ over the DB-eligible material:
\texttt{--gen-target N} drives it, \texttt{--gen-stop-frac} measures it as
$\mathrm{lag}/\mathrm{elig}$, and \texttt{--gen-giveup} writes off eligible gates
the DB genuinely cannot reach so the stop can fire. Read \texttt{lag=/elig=};
\texttt{wlag=} is the wide population.

Also required: \texttt{--no-local-verify}. The exhaustive per-rewrite check caps
support at 24 wires and \textbf{panics} on two wide gates. The whole-circuit
\texttt{--verify-every} check still runs.

\section{Options}

On \texttt{sss --cnot --gadgetize}; \texttt{gen\_sandwich\_gadget} reads the same
settings from the upper-case environment variables. With both counts at 0 the
encoding is off and the gadget is bit-identical to the linear build.

\begin{center}
\begin{tabular}{llll}
\toprule
flag & env var & default & meaning \\
\midrule
\texttt{--prod-k} & \texttt{PROD\_K} & 0 & base mask terms per value (statistical) \\
\texttt{--prod-deg} & \texttt{PROD\_DEG} & 2 & literals per base term, min 2 (algebraic) \\
\texttt{--prod-k-hi} & \texttt{PROD\_K\_HI} & 0 & tower terms per value \\
\texttt{--prod-deg-hi} & \texttt{PROD\_DEG\_HI} & 3 & tower degree; two deg-3 kill the deg-2 diagonal \\
\texttt{--prod-band} & \texttt{PROD\_BAND} & 0 (auto) & source-band width in extra wires \\
\texttt{--prod-rsrc} & \texttt{PROD\_RSRC} & 1 & mask re-source moves per inter-SG gap \\
\texttt{--prod-max-width} & \texttt{PROD\_MAX\_WIDTH} & 0 (off) & ladder to this width; \texttt{2} = all DB-eligible. ${\sim}6.4\times$ gates \\
\texttt{--prod-fill-nl} & \texttt{PROD\_FILL\_NL} & 0 (linear) & nonlinear cascaded band fill. Effectively free \\
\bottomrule
\end{tabular}
\end{center}

Mutually exclusive with the deferred-mask \texttt{--mask-cov} (RG4).
Auto band width is ${\approx}\max(\sqrt{4nk_{\mathrm{total}}},\,6,\,\deg_{\max}+3)$.

\textbf{Production:}
\begin{verbatim}
sss --cnot --gadgetize \
    --prod-k 1 --prod-deg 2 --prod-k-hi 2 --prod-deg-hi 3 --prod-fill-nl 2
\end{verbatim}

\section{Known gaps}

\begin{itemize}
\item \textbf{The statistical readout does not exist.} The exact-span ridge
measure saturates at $H=0.5$ for any $k\ge1$, so it cannot see what the term
count buys. Every $k$ decision is currently made against a measure blind to
precisely that axis. Highest-priority instrument.
\item \textbf{Band wires are body-static}---never written between the ports,
heavily read, hence trivially identifiable. That is what a restriction adversary
needs to condition on one. The mirror fill does not fix it; a rolling band would.
\item \textbf{Static distinguishers are not instrumented at all.} Every heatmap
here is execution-based. Gate-statistics along the circuit are unexamined, and
the fold-order fix of \S4 was found by reading code, not by a measurement that
would have caught it.
\item \textbf{The degree-4 tower} needs scratch decomposition of width-8
fragments.
\item \textbf{The DB match-rate cost} of width-6 gates is predicted by the
$(m,k,c)$ grid law but never measured directly.
\end{itemize}

\end{document}
